\documentclass[conference]{IEEEtran}

\usepackage{amsmath}
\usepackage{cite}
\usepackage{url}
\usepackage{array}

\begin{document}

\title{Milestone 2: Scientific Contextualization and Tradeoff Analysis\\
Loop Unrolling and Controller Delay in High-Level Synthesis}

\author{\IEEEauthorblockN{Christian Percival}
\IEEEauthorblockA{Hochschule Hamm-Lippstadt\\
Electronic Engineering}}

\maketitle

\section{Milestone 2: Scientific Contextualization and Tradeoff Analysis}

\subsection{Scientific Context}

\begin{itemize}
    \item Loop unrolling is a well-known optimization in compiler design.
    \item In traditional software compilation, it is mainly used to:
    \begin{itemize}
        \item reduce loop overhead,
        \item expose instruction-level parallelism,
        \item improve scheduling opportunities,
        \item increase performance on processors.
    \end{itemize}

    \item In high-level synthesis, the same optimization has different consequences.
    \item The unrolled loop body is not only compiled into more instructions, but can become more hardware structure.
    \item This means that loop unrolling affects:
    \begin{itemize}
        \item latency,
        \item datapath size,
        \item controller complexity,
        \item multiplexer complexity,
        \item resource usage,
        \item timing closure.
    \end{itemize}

    \item The reference paper is positioned between compiler optimization and hardware synthesis.
    \item Its main contribution is that it studies loop unrolling from the perspective of controller delay in HLS.
\end{itemize}

\subsection{Related Approaches}

\begin{itemize}
    \item Classical loop unrolling:
    \begin{itemize}
        \item Focuses on reducing loop overhead and improving software performance.
        \item Usually considers processor execution, instruction scheduling, and memory behavior.
        \item Does not directly consider FSM delay because software does not generate a hardware controller in the same way.
    \end{itemize}

    \item High-level synthesis scheduling:
    \begin{itemize}
        \item Decides which operations are executed in which clock cycles.
        \item Determines the number of control steps.
        \item Affects both latency and controller complexity.
    \end{itemize}

    \item Loop pipelining:
    \begin{itemize}
        \item Allows new loop iterations to start before previous iterations are complete.
        \item Often improves throughput.
        \item Can interact with loop unrolling.
        \item May require additional hardware resources and careful dependency handling.
    \end{itemize}

    \item Resource sharing:
    \begin{itemize}
        \item Reuses functional units such as adders and multipliers across multiple operations.
        \item Reduces area.
        \item Can increase multiplexer and controller complexity.
    \end{itemize}

    \item Exhaustive design-space exploration:
    \begin{itemize}
        \item Synthesizes many possible design versions.
        \item Can find good solutions.
        \item Becomes expensive when many loops and unroll factors are involved.
    \end{itemize}

    \item Estimation-based design-space exploration:
    \begin{itemize}
        \item Uses high-level models to estimate performance or delay.
        \item Avoids synthesizing every possible design.
        \item Is faster but less exact than full synthesis.
    \end{itemize}
\end{itemize}

\subsection{Main Tradeoff}

\begin{itemize}
    \item The central tradeoff is between latency reduction and controller-delay increase.
    \item Small unroll factors:
    \begin{itemize}
        \item keep the controller simpler,
        \item use fewer hardware resources,
        \item may still have relatively high latency.
    \end{itemize}

    \item Medium unroll factors:
    \begin{itemize}
        \item often give a good latency improvement,
        \item usually avoid the worst timing and area penalties,
        \item may be the most practical choice.
    \end{itemize}

    \item Large or full unroll factors:
    \begin{itemize}
        \item expose maximum parallelism,
        \item reduce the number of loop iterations,
        \item can greatly increase hardware resources,
        \item can increase FSM and multiplexer complexity,
        \item can create timing problems.
    \end{itemize}
\end{itemize}

\subsection{Latency Versus Clock Period}

\begin{itemize}
    \item A design with fewer cycles is not automatically faster.
    \item Total execution time depends on both cycle count and clock period.
    \item Loop unrolling can reduce the number of cycles.
    \item However, if unrolling increases the critical path, the clock period may need to become longer.
    \item Therefore, the final performance depends on:
\end{itemize}

\begin{equation}
Execution\ Time = Latency\ in\ Cycles \times Clock\ Period
\end{equation}

\begin{itemize}
    \item This means that a design with lower latency in cycles may still be worse if it has a much slower clock.
\end{itemize}

\subsection{Resource Tradeoffs}

\begin{itemize}
    \item Loop unrolling can increase resource usage because more operations are present in the loop body.
    \item More arithmetic operations may require more:
    \begin{itemize}
        \item adders,
        \item multipliers,
        \item DSP blocks,
        \item registers,
        \item multiplexers,
        \item memory ports.
    \end{itemize}

    \item If enough resources are available, unrolling can expose parallelism.
    \item If resources are limited, unrolled operations may still need to share the same hardware units.
    \item Resource sharing can reduce area but increase controller and multiplexer complexity.
\end{itemize}

\subsection{Timing Tradeoffs}

\begin{itemize}
    \item Unrolling can make the controller more complex.
    \item A more complex controller may have longer combinational logic paths.
    \item Larger multiplexers can also increase datapath delay.
    \item These effects can reduce the maximum clock frequency.
    \item Timing closure becomes more difficult when the unroll factor is too large.
\end{itemize}

\subsection{Scalability Tradeoffs}

\begin{itemize}
    \item The design-space exploration problem grows when there are multiple loops.
    \item Each loop can have several possible unroll factors.
    \item Exhaustively trying every combination can become impractical.
    \item The reference paper addresses this by:
    \begin{itemize}
        \item pruning the search space,
        \item considering promising unroll factors,
        \item using a priority function,
        \item estimating delay instead of synthesizing every candidate.
    \end{itemize}
\end{itemize}

\subsection{Comparison of Design Choices}

\begin{table}[h]
\centering
\caption{Tradeoff summary for loop unrolling in HLS}
\begin{tabular}{|p{0.22\linewidth}|p{0.32\linewidth}|p{0.32\linewidth}|}
\hline
\textbf{Design Choice} & \textbf{Advantages} & \textbf{Disadvantages} \\
\hline
No unrolling & Simple controller and low resource usage & Higher loop latency \\
\hline
Partial unrolling & Good balance between latency and complexity & Requires choosing a suitable factor \\
\hline
Full unrolling & Maximum exposed parallelism & High resource usage and possible timing problems \\
\hline
Exhaustive exploration & Can find strong design candidates & Very high synthesis time \\
\hline
Estimation-based exploration & Faster and more scalable & Less exact than full synthesis \\
\hline
\end{tabular}
\end{table}

\subsection{Relevance for Embedded Systems}

\begin{itemize}
    \item Embedded systems often have strict timing, area, power, and cost constraints.
    \item HLS is useful for embedded hardware accelerators, but optimization choices must be hardware-aware.
    \item Loop unrolling can improve performance, but may increase area and timing pressure.
    \item A balanced unroll factor is important for practical embedded hardware.
    \item The paper is relevant because it shows that controller delay must be considered, not only datapath latency.
\end{itemize}

\subsection{References and Related Work Situation}

\begin{itemize}
    \item Related work exists in compiler optimization, loop unrolling, software pipelining, HLS scheduling, timing estimation, and design-space exploration.
    \item However, the specific focus on controller delay caused by loop unrolling in HLS is more specialized.
    \item Therefore, not every loop-unrolling paper is equally relevant.
    \item The strongest related work should be selected from:
    \begin{itemize}
        \item loop unrolling in compilers,
        \item HLS timing estimation,
        \item HLS scheduling,
        \item loop pipelining,
        \item design-space exploration.
    \end{itemize}
\end{itemize}

\begin{thebibliography}{00}

\bibitem{kurra2007}
S. Kurra, N. K. Singh, and P. R. Panda, ``The Impact of Loop Unrolling on Controller Delay in High Level Synthesis,'' DATE, 2007.

\end{thebibliography}

\end{document}
